% main.tex – LaTeX version for Overleaf
% -------------------------------------------------
% This file contains the same content as the markdown version, but fully in LaTeX
% so Overleaf can compile it without needing the markdown+latex engine.

\documentclass[12pt]{article}
\usepackage[spanish]{babel}
\usepackage{amsmath,amssymb}
\usepackage{geometry}
\usepackage{hyperref}
\usepackage{xcolor}
\usepackage{listings}
\geometry{margin=1in}

% Code listing style
\lstset{basicstyle=\ttfamily\small,breaklines=true,backgroundcolor=\color{gray!10},frame=single}

\title{Arquitectura Frontend ⇄ Backend y Modelado Físico}
\author{Quizmate Development Team}
\date{\today}

\begin{document}
\maketitle

\section*{Resumen Ejecutivo}
El objetivo del proyecto es ofrecer una experiencia educativa interactiva en la que el usuario lanza un proyectil (lanza, arco, etc.) y observa su trayectoria basada en principios de cinemática y trigonometría.  
Todas las decisiones críticas (cálculo de la parábola, detección de impacto, cálculo de daño y "knockback") se realizan en el \textbf{backend} mediante un servidor Express + Prisma.  
El frontend, construido con \textbf{React} y \textbf{Phaser}, se limita a capturar la entrada del usuario, enviarla al servidor y reproducir visualmente los datos devueltos.

\section*{Stack Tecnológico}
\begin{itemize}
  \item \textbf{Frontend}: React 19, Vite, TypeScript, Phaser 3.
  \item \textbf{Backend}: Node.js, Express, TypeScript, Prisma (PostgreSQL), Socket.IO.
  \item \textbf{Comunicación}: REST (\texttt{/api/game/shoot}) + WebSockets.
  \item \textbf{Autenticación}: JWT (Bearer token) en cabecera \texttt{Authorization}.
\end{itemize}

\section*{Flujo de Interacción}
\begin{enumerate}
  \item El usuario ingresa velocidad, ángulo y opcionalmente ubicación del golpe ("head" o "body").
  \item Al pulsar "Disparar", el frontend envía una petición POST a \texttt{/api/game/shoot} con el cuerpo JSON:
\begin{verbatim}
{ "gameId": "<id>", "velocity": 75.5, "angleDegrees": 42, "direction": 1, "weaponId": "spear", "hitLocation": "head" }
\end{verbatim}
  \item El backend valida el JWT mediante \texttt{authMiddleware} y verifica que el jugador pertenezca a la partida.
  \item Se calcula la \textbf{trayectoria parabólica} (ver sección \ref{sec:trajectory}).
  \item Se detecta proximidad al oponente, se determina si el golpe es exitoso y se calcula daño y knockback (ver sección \ref{sec:damage}).
  \item Se persisten los cambios en la base de datos y se registra el disparo.
  \item El servidor emite \texttt{game\_update} vía Socket.IO para sincronizar al rival.
  \item El frontend recibe la respuesta/evento y actualiza UI y animaciones con Phaser.
\end{enumerate}

\section{Cálculo de la Trayectoria}\label{sec:trajectory}
El algoritmo implementado en \texttt{src/server/games/index.ts} emplea la fórmula clásica de movimiento bajo gravedad constante:
\begin{align}
  x(t) &= x_{0} + v\cos(\theta)\,d\,t \\
  y(t) &= y_{0} + v\sin(\theta)\,t - \frac{1}{2}g\,t^{2} + w\,t
\end{align}
Donde:
\begin{itemize}
  \item $x_{0},y_{0}$: posición inicial del lanzador.
  \item $v$: velocidad inicial.
  \item $\theta$: ángulo (rad).
  \item $d$: dirección horizontal (1 o -1).
  \item $g$: gravedad (valor fijo $0.5$).
  \item $w$: componente de viento horizontal (actualmente $0$, extensible).
\end{itemize}
El recorrido se genera iterando cada $\Delta t = 0.5\,\text{s}$ hasta que $y$ supera la altura inicial (impacto con el suelo).

\section{Daño y Knockback}\label{sec:damage}
Se calcula la distancia euclidiana entre el punto de impacto y el oponente:
\[ d = \sqrt{(x_{impact}-x_{op})^{2} + (y_{impact}-y_{op})^{2}} \]
Si $d < 55$ px, el disparo se considera exitoso ($\textit{isHit}=true$).

\subsection*{Rangos de Daño}
\begin{lstlisting}[language=typescript]
let damage = 0;
if (isHit) {
  const hitLocation = (req.body as any).hitLocation ?? 'body';
  if (hitLocation === 'head') {
    damage = Math.floor(Math.random() * (30 - 22 + 1)) + 22; // 22‑30
  } else {
    damage = Math.floor(Math.random() * (20 - 14 + 1)) + 14; // 14‑20
  }
}
\end{lstlisting}
El knockback se calcula como:
\[ k = d_{direction}\times \operatorname{round}(damage\times1.1) \]
Y se aplica mediante Prisma (`positionX` incrementado).

\section{Comunicación en Tiempo Real}
\begin{lstlisting}[language=typescript]
socketServer?.to(gameId).emit('game_update', {
  shooterId,
  trajectory,
  impactPoint,
  isHit,
  damage,
  knockback,
  newTargetPositionX
});
\end{lstlisting}
Cliente (React/Phaser):
\begin{lstlisting}[language=javascript]
socket.on('game_update', data => {
  // Dibujar trayectoria con Phaser, aplicar daño/knockback
});
\end{lstlisting}

\section{Renderizado en Phaser}
\begin{lstlisting}[language=javascript]
const graphics = this.add.graphics();
graphics.lineStyle(2, 0xffffff);
graphics.beginPath();
trajectory.forEach((p,i)=>{ if(i===0) graphics.moveTo(p.x,p.y); else graphics.lineTo(p.x,p.y);});
graphics.strokePath();
\end{lstlisting}

\section{Extensibilidad del Viento}
Para habilitar viento, envíe el campo `wind` en la petición y modifique `calculateTrajectory` para incluir $w\,t$ en $x(t)$. No se necesita tocar la lógica de daño.

\section*{Conclusiones}
El proyecto sigue el principio de \textbf{backend autoritativo}: toda la física y lógica de juego se calculan en el servidor, mientras el cliente solo renderiza. Esta arquitectura permite añadir nuevas variables (viento, gravedad variable) con mínimos cambios en el frontend.

\end{document}
